\documentclass{article}
\usepackage[utf8]{inputenc}
\usepackage{amsmath,amssymb}
\usepackage[most]{tcolorbox}
\usepackage{geometry}
\geometry{margin=2.5cm}

\newcommand{\step}[1]{\par\noindent\hangindent=3.5em\hangafter=1%
  \makebox[3.5em][l]{\textbf{#1.}}}
\newcommand{\steptag}[1]{\hfill{\small\textsf{[#1]}}}

\newtcolorbox{proofbox}[1]{
  colback=gray!5, colframe=gray!60,
  fonttitle=\bfseries, title={#1},
  breakable, enhanced,
  left=4pt, right=4pt, top=4pt, bottom=4pt,
  before skip=10pt, after skip=10pt
}

\begin{document}

\begin{proofbox}{Smooth polar-inverse upgrade: for every finite-dimensiona...}
\step{1} Smooth polar-inverse upgrade: for every finite-dimensional exact-unit epsilon\_r-C*-algebra and every delta > 0, if lem-stage1-smooth-unitary-atlas gives the smooth embedded atlas and Pi\_delta: calU x B\textasciicircum{}\{calH\}\_delta(J) -> S\_delta is the bijective C\textasciicircum{}1 local diffeomorphism of lem-stage1-polar-retraction onto an open set, then the same ambient-bilinear Pi\_delta is a smooth diffeomorphism and its same set-theoretic inverse (u\_delta, h\_delta) is smooth; no point or first derivative is changed. \steptag{validated} \\[6pt]
\step{1.1} Under the antecedent of node 1, lem-stage1-smooth-unitary-atlas supplies on calU the same smooth embedded graph atlas as its C\textasciicircum{}1 atlas, with no point or first derivative changed; B\textasciicircum{}\{calH\}\_delta(J) is an open subset of the finite-dimensional real vector space calH, so calU x B\textasciicircum{}\{calH\}\_delta(J) has the product smooth structure, while S\_delta is an open subset of the finite-dimensional ambient space by lem-stage1-polar-retraction and has its open-subset smooth structure. \steptag{validated} \\[6pt]
\step{1.1.1} By lem-stage1-smooth-unitary-atlas, under its stated hypotheses the existing C\textasciicircum{}1 graph charts cover calU and the same graph functions and charts form a smooth embedded-manifold atlas, without changing any point or first derivative. \steptag{validated} \\[4pt]
\step{1.1.2} By def-approximate-unitary-space, calH is the Hermitian real linear subspace of the finite-dimensional ambient complex vector space; hence calH is finite-dimensional over R and B\textasciicircum{}\{calH\}\_delta(J) is an open subset of calH with its standard smooth structure. \steptag{validated} \\[4pt]
\step{1.1.3} By lem-stage1-polar-retraction, S\_delta is open in the finite-dimensional ambient space, hence has the standard open-subset smooth structure. \steptag{validated} \\[4pt]
\step{1.1.4} The standard product construction applied to the smooth manifold calU and the open manifold B\textasciicircum{}\{calH\}\_delta(J) gives the asserted product smooth structure on calU x B\textasciicircum{}\{calH\}\_delta(J); together with the preceding open-subset structure on S\_delta this proves node 1.1. \steptag{validated} \\[4pt]
\step{1.2} Under the antecedent of node 1, the unchanged map Pi\_delta(U,H)=U bold-dot H is smooth from calU x B\textasciicircum{}\{calH\}\_delta(J) to S\_delta. \steptag{validated} \\[6pt]
\step{1.2.1} By def-epsilon-cstar-algebra, the ambient multiplication m(X,Y)=X bold-dot Y is bilinear. After choosing real bases in the finite-dimensional source and target spaces, each coordinate of m has the form sum\_\{i,j\} c\_\{kij\} x\_i y\_j, a polynomial; therefore m is C\textasciicircum{}infinity as a map of the underlying finite-dimensional real vector spaces. \steptag{validated} \\[4pt]
\step{1.2.2} Take any product chart on calU x B\textasciicircum{}\{calH\}\_delta(J) furnished by the unchanged smooth embedded atlas from lem-stage1-smooth-unitary-atlas and the identity chart on the ball, and use the ambient identity chart on the open set S\_delta. In these coordinates Pi\_delta is (x,H) maps to m(iota(psi\textasciicircum{}\{-1\}(x)),H), where iota composed with psi\textasciicircum{}\{-1\} is the smooth ambient parametrization of the embedded-manifold chart; the finite-coordinate polynomial formula from the preceding step and closure of smooth maps under composition show this coordinate representative is smooth. Since such charts cover the domain and Pi\_delta has values in S\_delta by lem-stage1-polar-retraction, Pi\_delta is smooth. \steptag{validated} \\[4pt]
\step{1.3} Under the antecedent of node 1, D Pi\_delta is a linear isomorphism at every point: this follows from the C\textasciicircum{}1 local-diffeomorphism assertion of lem-stage1-polar-retraction by differentiating a local C\textasciicircum{}1 inverse identity in charts. \steptag{validated} \\[6pt]
\step{1.3.1} Fix a point p of calU x B\textasciicircum{}\{calH\}\_delta(J). By the C\textasciicircum{}1 local-diffeomorphism assertion in lem-stage1-polar-retraction, there are neighborhoods of p and Pi\_delta(p) on which Pi\_delta has a C\textasciicircum{}1 inverse G, so in local Euclidean charts G composed with Pi\_delta and Pi\_delta composed with G are identity maps. \steptag{validated} \\[4pt]
\step{1.3.2} Applying the ordinary chain rule at p and Pi\_delta(p) to those two identities gives DG\_\{Pi\_delta(p)\} composed with DPi\_delta\_p = I and DPi\_delta\_p composed with DG\_\{Pi\_delta(p)\} = I. Thus DPi\_delta\_p is a linear isomorphism; since p was arbitrary, this holds everywhere. \steptag{validated} \\[4pt]
\step{1.4} Under the antecedent of node 1, the bijection Pi\_delta has a smooth inverse, and that inverse is exactly the already supplied set-theoretic pair (u\_delta,h\_delta). \steptag{validated} \\[6pt]
\step{1.4.1} Fix p in calU x B\textasciicircum{}\{calH\}\_delta(J), put q=Pi\_delta(p), and choose smooth charts alpha:W -> U subset R\textasciicircum{}n and beta:Z -> V subset R\textasciicircum{}n about p and q, shrinking W so Pi\_delta(W) subset Z. By node 1.2 the coordinate map F=beta composed with Pi\_delta composed with alpha\textasciicircum{}\{-1\} is smooth, and by node 1.3 its derivative at alpha(p) is invertible. \steptag{validated} \\[4pt]
\step{1.4.2} Apply GT-lee-2ed-thm-C.34 to F at alpha(p). It yields open neighborhoods U\_0 of alpha(p) and V\_0 of beta(q) on which F is a smooth diffeomorphism; translating through alpha and beta, Pi\_delta therefore has a smooth local inverse near q. \steptag{validated} \\[4pt]
\step{1.4.3} By the bijectivity assertion of lem-stage1-polar-retraction, the global set-theoretic inverse G=Pi\_delta\textasciicircum{}\{-1\} exists. On every local target neighborhood obtained in the preceding step, G equals the unique smooth local inverse there. These neighborhoods cover S\_delta as p ranges over the domain, so G is smooth because smoothness is local; this is exactly the local-gluing argument recorded in GT-lee-2ed-cor-C.36(b). \steptag{validated} \\[4pt]
\step{1.4.4} The inverse furnished by lem-stage1-polar-retraction is the pair (u\_delta,h\_delta). Since a bijection has only one set-theoretic inverse, the smooth global inverse G just obtained is exactly (u\_delta,h\_delta), proving node 1.4. \steptag{validated} \\[4pt]
\step{1.5} Under the antecedent of node 1, the smooth upgrade changes no point and no first derivative: lem-stage1-smooth-unitary-atlas retains the same graph charts through first order, Pi\_delta remains the identical ambient-bilinear map, and its inverse remains the identical C\textasciicircum{}1 set map (u\_delta,h\_delta), so uniqueness of derivatives gives the same first derivatives. \steptag{validated} \\[6pt]
\step{1.5.1} The external lemma lem-stage1-smooth-unitary-atlas expressly uses the same graph functions and graph charts as the C\textasciicircum{}1 atlas and changes no graph point or first derivative, so the source manifold upgrade is identity-on-data through first order. \steptag{validated} \\[4pt]
\step{1.5.2} The map in both lem-stage1-polar-retraction and the present conclusion is literally Pi\_delta(U,H)=U bold-dot H on the same domain and codomain, and the inverse in both is the unique set-theoretic inverse (u\_delta,h\_delta); thus none of these functions or their point values is replaced. \steptag{validated} \\[4pt]
\step{1.5.3} Before the upgrade these identical maps are C\textasciicircum{}1 by lem-stage1-polar-retraction. The derivative at a point is determined uniquely by the underlying C\textasciicircum{}1 function and the unchanged first-order charts, so merely proving that these same functions are C\textasciicircum{}infinity cannot change any first derivative. Together with the preceding two steps this proves node 1.5. \steptag{validated} \\[4pt]
\end{proofbox}

\end{document}
